\chapter{Gọi Cả Lớp Cùng Làm}
\chapterintro{Không để chỉ vài em gánh tiết học}{Tiết học dễ trôi khi chỉ có một nhóm nhỏ đang hoạt động thật.}

\section{Mỗi dãy một ý}
\begin{tinhhuong}{Tình huống}
Lớp chỉ có vài cánh tay quen thuộc, còn phần lớn ngồi im chờ qua tiết.
\end{tinhhuong}
\begin{goiydungngay}
Chia lớp theo dãy: mỗi dãy phải đóng góp đúng một ý, không cần dài.
\begin{itemize}
\item Gọi theo dãy giúp áp lực nhẹ hơn gọi cá nhân.
\item Ai cũng cảm thấy phần mình có trách nhiệm.
\item Hợp với cả lớp trầm và lớp đông.
\end{itemize}
\end{goiydungngay}
\begin{cauchot}
Khi trách nhiệm được chia đều, nhịp lớp thường sống lại.
\end{cauchot}
\ngancach

\section{Một câu cho mỗi bàn}
\begin{tinhhuong}{Tình huống}
Lớp bắt đầu rời rạc vì hoạt động đang quá cá nhân, thiếu cảm giác cùng nhau làm.
\end{tinhhuong}
\begin{goiydungngay}
Ra nhiệm vụ: mỗi bàn nói đúng một câu duy nhất về nội dung vừa học.
\begin{itemize}
\item Câu ngắn giúp lớp không ngại.
\item Nói theo bàn tạo cảm giác đồng đội.
\item Sau 5 bàn là lớp đã có lại năng lượng chung.
\end{itemize}
\end{goiydungngay}
\begin{cauchot}
Muốn lớp cùng học, hãy tạo những nhiệm vụ đủ nhỏ để ai cũng vào được.
\end{cauchot}
\ngancach

\section{Ai chưa nói sẽ nói trước}
\begin{tinhhuong}{Tình huống}
Đầu giữa tiết đã thấy vài em hoạt động quá nhiều, vài em gần như biến mất khỏi bài.
\end{tinhhuong}
\begin{goiydungngay}
Nhẹ nhàng thông báo: \textit{``Vòng này ưu tiên những bạn nãy giờ chưa lên tiếng.''}
\begin{itemize}
\item Không nêu tên ngay nếu lớp ngại.
\item Có thể cho các em chọn nói ngắn.
\item Không để lớp bị dẫn mãi bởi cùng một nhóm.
\end{itemize}
\end{goiydungngay}
\begin{cauchot}
Tiết học công bằng hơn thì cũng đỡ trôi hơn.
\end{cauchot}
\ngancach

\section{Mỗi góc lớp giữ một việc nhỏ}
\begin{tinhhuong}{Tình huống}
Hoạt động chung đang loãng vì vai trò của các nhóm chưa rõ.
\end{tinhhuong}
\begin{goiydungngay}
Chia lớp thành các khu nhỏ: góc này tìm ví dụ, góc kia chốt từ khóa, góc khác đặt câu hỏi.
\begin{itemize}
\item Vai trò rõ thì học sinh dễ nhập hơn.
\item Việc nhỏ giúp lớp không ngại.
\item Khi ghép lại, cả lớp có cảm giác thật sự cùng làm.
\end{itemize}
\end{goiydungngay}
\begin{cauchot}
Muốn cả lớp cùng chạy, đừng bắt mọi người làm đúng cùng một việc.
\end{cauchot}
\ngancach

\section{Giơ tay đồng loạt trước khi mời phát biểu}
\begin{tinhhuong}{Tình huống}
Nếu gọi phát biểu ngay, lớp lại rơi vào vài gương mặt quen thuộc.
\end{tinhhuong}
\begin{goiydungngay}
Hỏi một câu ngắn và yêu cầu cả lớp giơ tay chọn đáp án hoặc mức độ đồng ý trước.
\begin{itemize}
\item Cả lớp được nhập cuộc cùng lúc.
\item Sau đó mới mời 2 em giải thích.
\item Cách này rất hợp để tăng tỷ lệ tham gia nhanh.
\end{itemize}
\end{goiydungngay}
\begin{cauchot}
Cho lớp tham gia đồng loạt trước là cách mở cửa cho nhiều tiếng nói hơn.
\end{cauchot}
